% Part:sets-functions-relations
% Chapter: size-of-sets
% Section: introduction

\documentclass[../../../include/open-logic-section]{subfiles}

\begin{document}

\olfileid{sfr}{siz}{int}

\olsection{Pambuka}

Nalika Georg Cantor ngrembakakake teori himpunan ing taun 1870-an,
salah siji ancasane yaiku ndadekake gagasan kumpulan tanpa wates
bisa ditampa; wong ing Abad Tengahan bakal nyebut iki tanpa wates kang nyata.
Perangan wigati saka upaya iki yaiku andharane ngenani \emph{gedhene}
himpunan kang beda-beda. Yen $a$, $b$, lan $c$ kabeh beda,
mula miturut pangerten lumrah, himpunan $\{a, b, c\}$ \emph{luwih gedhe}
tinimbang $\{a, b\}$. Nanging kepiye himpunan tanpa wates?
Apa kabeh padha gedhene? Pranyata ora.

Gagasan wigati kang kapisan ing kene yaiku enumerasi.
Kita bisa ndhaptar saben himpunan winates kanthi ndhaptar kabeh !!{element}s.
Kanggo sawetara himpunan tanpa wates, kabeh !!{element}s uga bisa
didhaptar yen dhaptare dhewe oleh tanpa wates.
Himpunan kaya mangkono diarani !!{enumerable}.
Asil Cantor kang nggumunake, kang bakal kita mangerteni kanthi jangkep
ing pungkasan bab iki, yaiku ana himpunan tanpa wates kang ora !!{enumerable}.

\end{document}
